\section*{अध्याय \ref{ch:b3:virology} --- विषाणु-विज्ञान}

\begin{solution}{exo:b3:virology:1}
पोलियो: IV, धन-लड़ी RNA --- कोई RNA-निर्भर RNA पॉलीमरेज़।
इन्फ़्लुएंज़ा: V, ऋण-लड़ी RNA --- अपनी ही RNA पॉलीमरेज़, जो \omterm{def:b3:virology:virus}{विषाणुकण} में
ढोई जाती है। HIV: VI --- प्रतिलोम अनुलेखक (और समाकलक)। हर्पीज़
सिंप्लेक्स: I, द्विलड़ी DNA --- परपोषी सिद्धांततः सब कुछ दे सकता था, पर
वह अपनी DNA पॉलीमरेज़ और थायमिडीन काइनेज़ साथ लाता है।
हेपटाइटिस B: VII --- अपने पूर्वजीनोमी RNA को DNA में उतारने के लिए कोई
प्रतिलोम अनुलेखक। रोटावायरस: III, द्विलड़ी RNA --- कण के भीतर कोई
RNA-निर्भर RNA पॉलीमरेज़।
\end{solution}

\begin{solution}{exo:b3:virology:2}
फ़ेज का DNA ($^{32}$P) जीवाणुओं में घुसा और प्रोटीन
($^{35}$S) बाहर रह गया तथा झाड़ा जा सका; संतति फ़ेज को
$^{32}$P विरासत में मिला। इसलिए फ़ेज जो अंतःक्षेपित करता है और जो नए
फ़ेज बनवाता है वह DNA है। यदि निर्देश प्रोटीन ढोता, तो
$^{35}$S भीतर मिलता और संतति तक पहुँचता।
\end{solution}

\begin{solution}{exo:b3:virology:3}
$84/(0.1\times 10^{-7}) = 8.4\times 10^{9}$ पट्ट-निर्माणी इकाइयाँ प्रति mL।
सूक्ष्मदर्शी हर कण गिनता है, खाली \omterm{def:b3:virology:virus}{कैप्सिड}, दोषपूर्ण जीनोम वाले कण और
वे भी जो जुड़ने या अंतःक्षेपण में विफल रहते हैं; केवल वे कण पट्ट बनाते
हैं जो कोई संक्रमण पूरा करते हैं।
\end{solution}

\begin{solution}{exo:b3:virology:4}
जुड़ाव: प्रवेश-निरोधक (मैराविरॉक), उदासीनकारी प्रतिरक्षी।
प्रवेश/संलयन: संलयन-निरोधक, अंतःकायी अम्लीकरण रोकने वाली औषधें।
अनावरण: कैप्सिड-बंधक औषधें (प्लेकोनैरिल; लेनाकापाविर)। अभिव्यक्ति
और प्रतिकृति: \omterm{def:b3:virology:drugs}{न्यूक्लिओसाइड अनुरूप}, प्रतिलोम-अनुलेखक तथा
समाकलक निरोधक, इंटरफ़ेरॉन, \omterm{def:b3:rna-regulation:rnai}{RNA व्यवधान}। संयोजन: \omterm{def:b3:virology:drugs}{प्रोटिएज़
निरोधक} (अविदलित बहुप्रोटीन जुड़ ही नहीं सकते)। मोचन:
न्यूरामिनिडेज़ निरोधक (इन्फ़्लुएंज़ा), और कोशिकाविषी T~कोशिकाएँ जो
मोचन से पहले कोशिका मार देती हैं।
\end{solution}

\begin{solution}{exo:b3:virology:5}
\qty{30}{kDa} का एक प्रोटीन $30\,000/110 = 273$ अवशेष, यानी $818$
न्यूक्लियोटाइड है: \qty{0.82}{kb}, यानी \qty{4.5}{kb} जीनोम का \qty{18}{\%}।
उन $180$ प्रतियों को अलग-अलग कूटबद्ध करने में \qty{147}{kb} लगते,
यानी पूरे जीनोम से तैंतीस गुना; स्वयं \omterm{def:b3:virology:virus}{कैप्सिड} प्रोटीन
(\qty{5.4}{MDa}) जीनोम (\qty{1.5}{MDa}) से चार गुना भारी है।
\end{solution}

\begin{solution}{exo:b3:virology:6}
$V = 10^{5}(3e^{-0.5t} - 0.5e^{-3t})/2.5$: $t = 0.5$: $8.9\times
10^{4}$; $t = 2$: $4.4\times 10^{4}$; $t = 10$: $8\times 10^{2}$। वह
$50$ से नीचे तब जाता है जब $1.2\times 10^{5}e^{-0.5t} = 50$: $t = 2\ln(2400) \approx
\qty{16}{d}$।
\end{solution}

\begin{solution}{exo:b3:virology:7}
$uL = 0.36$; त्रुटि-रहित अंश $e^{-0.36} = 0.70$। तिगुना: $uL =
1.08$, $e^{-1.08} = 0.34$ --- अब अधिकांश संतति उत्परिवर्तन ढोती है और
सबसे अनुकूल अनुक्रम देहली की ओर तनु होता जाता है। $u = 10^{-4}$ वाले
\qty{30}{kb} के किसी कोरोना \omterm{def:b3:virology:virus}{विषाणु} का $uL = 3$ होता और केवल \qty{5}{\%}
प्रतियाँ त्रुटि-रहित, यानी देहली के पार; उसकी एक्सोन्यूक्लिएज़ $u$ को
$10^{-5}$ पर ले आती है, $uL = 0.3$, यानी तीन-चौथाई त्रुटि-रहित।
\end{solution}

\begin{solution}{exo:b3:virology:8}
अपवाह: हीमैग्लुटिनिन और न्यूरामिनिडेज़ में बिंदु उत्परिवर्तन
प्रतिरक्षी-चयन के अंतर्गत जमा होते हैं, इसलिए हर मौसम की प्रजाति
कुछ नई होती है और पिछले वर्ष की प्रतिरक्षा कुछ पुरानी। विस्थापन: किसी
एक कोशिका का किसी मानव और किसी पशु प्रजाति से सह-संक्रमण आठों खंड
पुनर्मिश्रित कर देता है, और पक्षियों या सूअरों का कोई हीमैग्लुटिनिन
उपप्रकार मनुष्यों के अनुकूल किसी \omterm{def:b3:virology:virus}{विषाणु} में आ जाता है। नए उपप्रकार के
प्रति किसी के पास प्रतिरक्षी नहीं, इसलिए पूरी समष्टि एक साथ संवेद्य है
--- यही महामारी की शर्त है, जो 1918, 1957, 1968 और 2009 में हर बार किसी
नए हीमैग्लुटिनिन से पूरी हुई।
\end{solution}

\begin{solution}{exo:b3:virology:9}
हर्पीज़ \omterm{def:b3:virology:virus}{विषाणु} की थायमिडीन काइनेज़ ऐसिक्लोविर का फ़ॉस्फ़ोरिलन करती है,
जिसे कोशिकीय काइनेज़ मुश्किल से छूती हैं; फिर कोशिकीय काइनेज़ त्रिफ़ॉस्फ़ेट
पूरा करती हैं, जिसे विषाणुजन्य DNA पॉलीमरेज़ बैठा लेती है और, 3$'$
हाइड्रॉक्सिल न होने से, शृंखला समाप्त हो जाती है। असंक्रमित कोशिकाएँ
सक्रिय रूप कभी नहीं बनातीं, और विषाणुजन्य पॉलीमरेज़ उसे कोशिकीय
पॉलीमरेज़ से अधिक पसंद करती है। प्रतिरोध: विषाणुजन्य थायमिडीन काइनेज़
का ह्रास या बदलाव (सबसे आम), या विषाणुजन्य पॉलीमरेज़ के ऐसे
उत्परिवर्तन जो उस अनुरूप को बाहर रखें।
\end{solution}

\begin{solution}{exo:b3:virology:10}
$T$ के अचर रहने पर समीकरण $(I, V)$ में रैखिक हैं और उनका आव्यूह
$\begin{pmatrix} -\delta & \beta T \\ p & -c \end{pmatrix}$ है, जिसका अनुरेख
$-(\delta + c) < 0$ है और सारणिक $\delta c - \beta Tp$। अनुरेख
ऋणात्मक होने से कोई एक अभिलक्षणिक मान ठीक तभी धनात्मक होता है जब
सारणिक ऋणात्मक हो: $\beta Tp > c\delta$, यानी $R_{0} = \beta
Tp/(c\delta) > 1$। गुणक: $p/\delta$ वह संख्या है जितने \omterm{def:b3:virology:virus}{विषाणुकण} कोई
संक्रमित कोशिका अपने जीवन में बनाती है; $\beta T/c$ वह संख्या है जितनी
कोशिकाएँ कोई \omterm{def:b3:virology:virus}{विषाणुकण} अपने जीवन में संक्रमित करता है; उनका गुणनफल वह
संख्या है जितनी संक्रमित कोशिकाएँ एक संक्रमित कोशिका बनाती है।
\end{solution}

\begin{solution}{exo:b3:virology:11}
$10^{6} = 2^{20}$: बीस अर्ध-आयु, $880$ महीने, लगभग $73$ वर्ष ---
किसी जीवन से भी लंबा। प्रतिकृति रोकने वाली औषधें ऐसे प्रोविषाणु को
छूती ही नहीं जिसका अनुलेखन न हो रहा हो; भंडार बना रहता है, और जब
चिकित्सा रुके तो कोई एक फिर सक्रिय होती कोशिका संक्रमण दुबारा शुरू कर
देती है। किसी इलाज को भंडार हटाना या स्थायी रूप से चुप करना पड़ेगा:
प्रसुप्त कोशिकाएँ मारिए (चिकित्सा के अंतर्गत उन्हें फिर सक्रिय कीजिए
ताकि वे मरें या मारी जाएँ), प्रोविषाणु उच्छेदित या निष्क्रिय कीजिए, या
प्रतिरक्षा-तंत्र को ऐसे तंत्र से बदल दीजिए जिसे \omterm{def:b3:virology:virus}{विषाणु} संक्रमित न कर सके।
\end{solution}

\begin{solution}{exo:b3:virology:12}
$uL = 0.03$ उत्परिवर्तन प्रति जीनोम; प्रति दिन $10^{9}$ जीनोम $3\times
10^{7}$ उत्परिवर्तन ढोते हैं, जो $9\times 10^{4}$ संभव एकल
उत्परिवर्तनों में बँटते हैं: हर एक प्रति दिन लगभग $300$ बार उठता है।
कोई विशिष्ट दोहरा उत्परिवर्ती प्रति जीनोम
$10^{-12}$ से उठता है, यानी प्रति दिन $10^{-3}$ --- तीन वर्षों में एक बार; कोई
तिहरा $10^{-18}$ से, यानी कभी नहीं। स्वतंत्र प्रतिरोध-उत्परिवर्तनों
वाली दो औषधें सिद्धांततः पर्याप्त होतीं, और तीन खराब पालन के लिए
गुंजाइश देती हैं। एक ही एंज़ाइम से बँधती दो औषधें किसी एक ऐसे
उत्परिवर्तन से हार सकती हैं जो स्थल या एंज़ाइम की संरूपता बदल दे, और
उनके प्रतिरोध-उत्परिवर्तन स्वतंत्र नहीं हैं; वे दो से कम गिनी जाती हैं।
\end{solution}

\begin{solution}{pb:b3:virology:1}
\textbf{1.} $2\times 10^{5}\times 1.5\times 10^{4} = 3\times 10^{9}$
\omterm{def:b3:virology:virus}{विषाणुकण}।
\textbf{2.} प्रति दिन साफ़ $cV = 3\times 3\times 10^{9} \approx 10^{10}$;
उत्पादन भी उतना ही।
\textbf{3.} $I^{*} = cV^{*}/p = 9\times 10^{9}/500 = 1.8\times 10^{7}$
संक्रमित कोशिकाएँ; प्रति दिन मरती $\delta I^{*} = 9\times 10^{6}$।
\textbf{4.} RNA $2\times 9700\times 330 = 6.4\times 10^{6}$ Da; प्रोटीन
$2000\times 25\,000 = 5\times 10^{7}$ Da; कुल $5.6\times 10^{7}$ Da
$\approx 10^{-16}$ g। सारे \omterm{def:b3:virology:virus}{विषाणुकण} मिलकर: $3\times 10^{9}\times
10^{-16} = 3\times 10^{-7}$ g, यानी एक माइक्रोग्राम का तीसरा भाग ---
रेत के किसी दाने से तीन हज़ार गुना कम।
\textbf{5.} भंडार का $1.8\times 10^{7}/10^{12} = 2\times 10^{-5}$।
प्रति दिन $9\times 10^{6}$ मृत्युओं के दस वर्ष $3\times 10^{10}$
कोशिकाएँ हैं, यानी भंडार का तीन प्रतिशत: भंडार अंकगणितीय घटाव से नहीं
चुकता, बल्कि इसलिए कि दशक भर की बलात् पुनर्जननी और चिरकारी सक्रियता
उन्हें बदलने की क्षमता चुका देती है।
\textbf{6.} किसी भी क्षण $9\times 10^{6}/24 \approx 4\times 10^{5}$
मुर्दे साफ़ किए जा रहे हैं।
\textbf{7.} $V = 2\times 10^{5}(3e^{-0.5t} - 0.5e^{-3t})/2.5$: $t = 1$:
$1.4\times 10^{5}$; $t = 3$: $5.4\times 10^{4}$; $t = 7$: $7\times
10^{3}$; $t = 14$: $2\times 10^{2}$।
\textbf{8.} $2.4\times 10^{5}e^{-0.5t} = 50$: $t = 2\ln(4800) \approx 17$
दिन।
\textbf{9.} दिन 3--10 धीमी प्रावस्था पर पड़ते हैं: वहाँ $\ln V$ की ढाल
$-\delta$ है (सात दिनों में $5.4\times 10^{4}$ से $1.6\times 10^{3}$,
यानी $\delta = 0.50$)। तेज़ प्रावस्था केवल घंटों चलती है ($1/c$), इसलिए
$c$ के लिए पहले दिन हर कुछ घंटे प्रतिदर्श चाहिए, जिन्हें दो-चरघातांकी
सूत्र से फिट किया जाए।
\textbf{10.} प्रसुप्त रूप से संक्रमित कोशिकाएँ, जो फिर सक्रिय होते समय
\omterm{def:b3:virology:virus}{विषाणु} छोड़ती हैं; उनकी क्षय-दर प्रति महीना $\ln 2/44$ है, यानी प्रति
दिन लगभग $5\times 10^{-4}$।
\textbf{11.} \omterm{def:b3:virology:virus}{विषाणुकण} $\ln 2/3 = \qty{0.23}{d}$, लगभग \qty{5.5}{h};
संक्रमित कोशिका $\ln 2/0.5 = \qty{1.4}{d}$।
\textbf{12.} अचर भार ऐसा लगता था मानो \omterm{def:b3:virology:virus}{विषाणु} कुछ कर ही न रहा हो;
समीकरण उसे ऐसी स्थायी अवस्था दिखाते हैं जिसमें हर दिन $10^{10}$
\omterm{def:b3:virology:virus}{विषाणुकण} बनते और साफ़ होते हैं और $10^{7}$ T~कोशिकाएँ संक्रमित तथा मारी
जाती हैं।
\textbf{13.} प्रति जीनोम $9700\times 10^{-5} \approx 0.1$ उत्परिवर्तन;
प्रति दिन $10^{10}\times 0.1 = 10^{9}$ उत्परिवर्तन।
\textbf{14.} $3\times 9700 \approx \num{29000}$ संभव एकल उत्परिवर्तन;
हर एक प्रति दिन $10^{9}/\num{29000} \approx 3\times 10^{4}$ बार उठता है।
\textbf{15.} विशिष्ट दोहरा: प्रति जीनोम $10^{-10}$, अनुपचारित अवस्था
में लगभग प्रति दिन एक; विशिष्ट तिहरा: $10^{-15}$, प्रति दिन $10^{-5}$
--- तीन सदियों में एक बार।
\textbf{16.} प्रति दिन $10^{7}\times 10^{-15} = 10^{-8}$, प्रति वर्ष
$4\times 10^{-6}$।
\textbf{17.} एक औषध हट जाने पर बाकी दो के प्रति प्रतिरोधी उत्परिवर्ती
पर्याप्त हैं: कोई विशिष्ट दोहरा $10^{-10}$ पर; उस सप्ताह के $7\times 10^{7}$
जीनोम $7\times 10^{-3}$ प्रत्याशित देते हैं --- फिर भी असंभाव्य,
जब तक चूक के दौरान भार प्रति दिन $10^{10}$ की ओर न लौट आए, जब वह प्रति
दिन एक हो जाता है और बच निकलना संभाव्य। छूटी मात्राएँ ही वह रास्ता हैं
जिससे प्रतिरोध आता है।
\textbf{18.} ये अनुरूप प्रतिलोम अनुलेखक का सक्रिय स्थल साझा करते हैं,
और वहाँ का एक उत्परिवर्तन (या थायमिडीन-अनुरूप उत्परिवर्तनों का कोई
समुच्चय) एंज़ाइम के उनमें से कई को सँभालने का ढंग बदल देता है:
पार-प्रतिरोध। दो न्यूक्लिओसाइड दो स्वतंत्र औषधों से कम गिने जाते हैं,
इसलिए नियम-पद्धतियाँ वर्ग मिलाती हैं --- कोई न्यूक्लिओसाइड, कोई समाकलक
निरोधक, कोई \omterm{def:b3:virology:drugs}{प्रोटिएज़ निरोधक} या कोई अ-न्यूक्लिओसाइड।
\textbf{19.} $20$ अर्ध-आयु, $880$ महीने, $73$ वर्ष; $10^{4}$ तक,
$\log_{2}100 = 6.6$ अर्ध-आयु, $290$ महीने, $24$ वर्ष।
\textbf{20.} एक \omterm{def:b3:virology:virus}{विषाणुकण} से $3\times 10^{9}$ तक प्रति पीढ़ी $R_{0} = 8$
पर: $\ln(3\times 10^{9})/\ln 8 = 10.5$ पीढ़ियाँ, यानी तीन सप्ताह।
\textbf{21.} आघात और वध: चिकित्सा के अंतर्गत प्रसुप्त कोशिकाएँ फिर
सक्रिय कीजिए ताकि वे मरें या मारी जाएँ --- अड़चन: कोई कारक उन सबको फिर
सक्रिय नहीं करता, और फिर सक्रिय हुई कोशिकाएँ भरोसे से मारी नहीं जातीं।
जीन-संपादन: प्रोविषाणु उच्छेदित कीजिए या रोगी की कोशिकाओं में ग्राही
CCR5 नष्ट कीजिए --- अड़चन: हर कोशिका तक पहुँचना। (मज्जा को
CCR5-अपर्याप्त दाता कोशिकाओं से बदलने ने मुट्ठी भर रोगियों को ठीक किया
है, ऐसे जोखिम पर जो तभी स्वीकार्य है जब प्रत्यारोपण वैसे भी ज़रूरी हो।)
\textbf{22.} $1 - 1/R_{0}$: $R_{0} = 3$ के लिए \qty{67}{\%}, $5$ के लिए
\qty{80}{\%}।
\textbf{23.} चिकित्सा रोगी के भीतर $\beta \to 0$ कर देती है, भार प्रमेय
से लगभग शून्य तक गिर जाता है, और संचरण, जो भार के समानुपाती है, रुक
जाता है; हर संक्रमित व्यक्ति का उपचार समष्टि के $R_{0}$ को एक से नीचे
ले आता है।
\textbf{24.} $f = (1 - 1/R_{0})/E$: $R_{0} = 4$ और $E = 0.7$ के साथ $f =
0.75/0.7 = 1.07$ --- सबसे भी अधिक: अकेला \omterm{def:b3:virology:drugs}{टीका} देहली तक नहीं पहुँच
सकता; $E = 0.9$ के साथ $f = 0.83$।
\textbf{25.} प्रति दिन लगभग $10^{10}$ \omterm{def:b3:virology:virus}{विषाणुकण} बनते हैं; लगभग दिन
$17$ को पकड़ से बाहर; चिकित्सा पर प्रति वर्ष कोई $4\times 10^{-6}$
तिहरे उत्परिवर्ती।
\end{solution}
